\chapter{Характерстики сигнала}

\section{Описание сигналов в частотной области}

Комплексный спектр
\begin{align}
	P(if) = \int_{-\infty}^{\infty} s(t)e^{-j2\pi ft} dt = F(s(t))
\end{align}

Сигнал
\begin{align}
	st(t) = \int_{-\infty}^{\infty} P(if))e^{-j2\pi ft} df = F^{-1}(P(if))
\end{align}

Амплитудный спектр
\begin{align}
	P(f) = |P(if)| = \sqrt{Re^{2}P(if) + Im^{2}P(if)}
\end{align}

Фазовый спектр
\begin{align}
	\Phi(f) = arg(P(if)) = \arctan{\frac{ImP(if)}{ReP(if)}}
\end{align}

Спектральная плотность энергии:
\begin{align}
	G(f) = |P(if)|^{2} = P(if)\overline{P}(if) \equiv P^{2}(if)
\end{align}

\section{Описание сигналов во временной области}

Корреляционная функция 
\begin{align}
	B_{s}(\tau) = \int_{-\infty}^{\infty} s(t)s(t - \tau) dt \equiv \int_{-\infty}^{\infty} s(t + \tau)s(t) dt
\end{align}

Взаимная корреляционная функция
\begin{align}
	B_{s, u}(\tau) = \int_{-\infty}^{\infty} s(t)s(t - \tau)u dt
\end{align}

Коэффициент корреляции:
\begin{align}
	R(\tau) = \frac{B(\tau)}{B(0)}
\end{align}

Теорема Винера-Хинчина
\begin{align}
	B(\tau) = F^{-1}(G(f)) = \int_{-\infty}^{\infty} G(f)e^{-j2\pi ft} df \\
	G(f) = F(B(\tau)) = \int_{-\infty}^{\infty} B(\tau)e^{-j2\pi ft} d\tau \nonumber
\end{align}

\chapter{Информационные характеристики сигналов}

\section{Эффективная ширина спектра сигнала}

Центральная частота сигнала
\begin{align}
	f_{oo} = \frac{\int_{f_{l}}^{f_{h}} f G(f) df}{\int_{f_{l}}^{f_{h}} G(f) df}
\end{align}

Эффективная ширина спектра для НЧ-сигнала
\begin{align}
	F_{e}: G(f) \leq \epsilon \cdot G(f_{oo})
\end{align}

Эффективная ширина спектра для полосового сигнала
\begin{align}
	F_{e}: G(f - f_{oo}) \leq \epsilon \cdot G(f_{oo})
\end{align}

Эффективная длительность сигнала
\begin{align}
	\tau_{e}: |B(\tau)| \leq \rho \cdot B(0)
\end{align}

База сигнала
\begin{align}
	B_{o} = F_{e} \cdot \tau_{e}
\end{align}

Размерность сигнала
\begin{align}
	n = 2 \lfloor B_{o} \rfloor = 2 \lfloor F_{e} \cdot \tau_{e} \rfloor = n_{dim}
\end{align}

\chapter{Дискретизация и квантование сигнала}

\section{Дискретизация и восстановление сигналов}

Максимальная частота дискретизации
\begin{align}
	f_{max} = F
\end{align}

Минимальная частота дискретизации
\begin{align}
	f_{d_min} = \frac{\tau_{e}}{\Delta_{d}} = 2F \text{ (для ЧО-сигналов)} \\
	f_{d_min} = 2F_{e} \text{ (для не ЧО-сигналов)}\nonumber
\end{align}
 
Количество отсчетов
\begin{align}
	N = \lfloor \frac{\tau_{e}}{\Delta_{d}} \rfloor = \lfloor \tau_{e}f_{d_min} \rfloor = \lfloor 2F_{e}\tau_{e} \rfloor \approx 2 \lfloor B_{o}\rfloor  = n_{dim}
\end{align}

\section{Квантование дискретно-непрерывных сигналов}

Равномерный шаг квантования
\begin{align}
	\Delta = \frac{b_{max} - b_{min}}{M - 1}
\end{align}

Шум квантования
\begin{align}
	\beta_{k} = b_{k} - b_{Qk}
\end{align}

Дисперсия шумов квантования
\begin{align}
	D_{\beta} = M[\beta_{k}^{2}] = \frac{\Delta^{2}}{12} 
\end{align}

Среднеквадратическое значение шумов квантования
\begin{align}
	\sigma_{\beta} = \sqrtsign{D_{\beta}} = \frac{\Delta}{2\sqrtsign{3}}
\end{align}

Отсчет цифрового сигнала
\begin{align}
	b_{Qk} = b_{k}^{(i)} = q_{k}^{(i)}\Delta,
\end{align}
где 
\begin{itemize}
	\item $q_{k}^{(i)} = p^{(i)} \cdot j_{k}$ -- 
	\item $j_{k} = |i_{k} - (L - 1)|$ -- значение интервала наблюдения
	\item $M = 2L - 1$ -- количество уровней
	\item $L$ -- количество состояний
	\item $p_{(i)}$ -- знак
\end{itemize}

\chapter{Кодирование цифровых сигналов}

Процедура десятичного кодирования цифрового сигнала (в целочисленный код)
\begin{align}
	q_{k}^{(i)} = p^{(i)}j_{k} = p^{(i)}(q_{i, n - 1}, ..., q_{i, 0}), q_{i, l} = \overline{0, 9}
\end{align}

Процедура двоичного кодирования цифрового сигнала (в двоичный код) (ИКМ-сигнал)
\begin{align}
	w_{k}^{(i)} = p^{(i)}(w_{i, n}, ..., w_{i, 0}), w_{i, l} = \overline{0, 1}
\end{align}

Алгоритм получения первых n старших разрядов является рекуррентным
\begin{align}
	w_{i, l-1} = Q(i) = \frac{j - \sum_{l=l_{0}}^{n-1} w_{i, l_{0}}2^{l}}{2^{i}},
\end{align}
где 
% \lfloor \rfloor
\begin{itemize}
	\item $w_{i, j - 1} = \lfloor \frac{j}{2^{n - 1}} \rfloor$;
	\item $j = i + L + 1$;
	\item $n = \lfloor \log_{2}(L - 1) \rfloor$.
\end{itemize}

Старший разряд, определяющий знак $p^{(i)}$ цифрового сигнала
\begin{align}
	w_{i, n} = 
	\begin{cases}
		1, p^{(i)} = 1\\
		0, p^{(i)} = -1
	\end{cases}
\end{align}

Оператор, обратный двоичному кодированию 
\begin{align}
	q_{k}^{(i)} = Q^{-1}(w_{i, k}) = p^{(i)}q_{k} = p^{(i)} \sum_{l=0}^{n - 1} 2^{l}w_{i, l},\\
	\text{где }
	p^{(i)} = 
	\begin{cases}
		\text{\hspace*{0.7em}}1, w_{i, n} = 1\\
		-1, w_{i, n} = 0
	\end{cases}\nonumber
\end{align}

Алгоритм m-ичного кодирования, при котором десятичные числа записываются в n-позиционной системе счисления по основанию m>2 в виде совокупности разрядов
\begin{align}
	v_{i, k} = (v_{i, n}, ..., v_{i, 0}), v_{i, l} = \overline{0, m - 1}
\end{align}
Алгоритмы m-ичного кодирования и декодирования строятся аналогично двоичному варианту.

\chapter{Энтропия дискретного источника информации}

Собственная случайная информация Шеннона ДИС
\begin{align} 
	\label{shannon_info}
	I(x) = I(x_1, x_2) = -\log_{2}{p_{1, 2}(x_1, x_2)}
\end{align}

Собственная энтропия Шеннона ДИС
\begin{align}
	\label{shannon_entropy}
	H(x) = M(I(x)) = -\sum_{k=1}^{L} \sum_{q=1}^{L} p_{1, 2}(a_k, a_q)\log_{2}{p_{1, 2}(a_k, a_q)}
\end{align}

Информация Хартли двухмерного ДИС $X = \{x_1, x_2\}$
\begin{align}
	\label{hart}
	I_{Hart}(x) = -\log_{2}{\frac{1}{L^2}} = 2\log_{2}{L}
\end{align}

Собственная условная энтропия Шеннона ДИС
\begin{align}
	\label{shannon_entropy_if}
	H(x_1 | x_2) =  -\sum_{k=1}^{L} \sum_{q=1}^{L} p_{1, 2}(a_k, a_q)\log_{2}{p_{1, 2}(a_k | a_q)}
\end{align}

Собственная взаимная энтропия Шеннона двух случайных величин одного ДИС
\begin{align}
	\label{shannon_entropy_double}
	H(x \wedge y) = \sum_{r=1}^{M} \sum_{k=1}^{L} p_(a_k, b_r)\log_{2}{\frac{p(a_k | b_r)}{p(a_k)}}
\end{align}

Для одномерного ДИС $X = \{x\}$ из одной случайной величины с одномерным ДВР 
$p(a_k) = P\{x_1 = a_k\}, a_k \in A$ выражения \ref{shannon_info}--\ref{hart} принимают вид:
\begin{gather}
	\label{single_dim}
	I(x) = -\log_{2}{p(x)}, \\\nonumber
	I_{Hart}(x) = -\log_{2}{\frac{1}{L}} = \log_{2}{L}, \\
	H(x) = M(I(x)) = \sum_{k=1}^{L} p_(a_k)\log_{2}{p_(a_k)}. \nonumber
\end{gather}

Для ДИС из двух случайных величин $\{x_1, x_2\} = x$ с заданным двухмерным ДРВ  $p_{1,2}(a_k, a_q)$ собственная случайная информация (\ref{shannon_info}) и энтропии (\ref{shannon_entropy}--\ref{hart}) каждой из этих случайных величин могут быть вычислены по формулам \ref{single_dim} после нахождения их одномерных распределений из условия согласованности ДРВ по формулам:
\begin{flalign}
	p(a_k) = \sum_{q=1}^{L} p(a_k, a_q), p(a_q) = \sum_{k=1}^{L} p(a_k, a_q). 
\end{flalign}

\section{Взаимная информация и энтропия двух ДИС}

Совместная взаимная энтропия двух ДИС
\begin{align}
	H(x, y) = M(I(x, y)) = -\sum_{r=1}^{M} \sum_{k=1}^{L} p(a_k, b_r)\log_{2}{p(a_k, b_r)}
\end{align}

Условная взаимная энтропия двух ДИС
\begin{align}
	H(x | y) = -\sum_{r=1}^{M} \sum_{k=1}^{L} p(a_k, b_r)\log_{2}{p(a_k | b_r)}
\end{align}

Взаимная энтропия двух ДИС
\begin{align}
	H(x \wedge y) = \sum_{r=1}^{M} \sum_{k=1}^{L} p(a_k, b_r)\log_{2}{\frac{p(a_k | b_r)}{p(a_k)}}
\end{align}

Собственные взаимные энтропии двух ДИС
\begin{gather}
	H(x) = -\sum_{k=1}^{L} p(a_k)\log_{2}{p(a_k)}, \\
	H(y) = -\sum_{r=1}^{M} p(b_r)\log_{2}{p(b_r)} \nonumber
\end{gather}

\section{Основные тождества и неравенства для энтропий ДИС}

\begin{align}
	--
\end{align}

\section{Собственные энтропии ДИС}

Энтропия ДИС-БП
\begin{align}
	--
\end{align}

Энтропия ДИС-CM
\begin{align}
	--
\end{align}

\section{Экстремальные свойства энтропии}

\section{Удельная энтропия, производительность, насыщенность и избыточность ДИИ}

Удельная энтропия
\begin{align}
	--
\end{align}

Удельная эргодическая энтропия
\begin{align}
	H_{\infty}(x) = \lim_{n \rightarrow \infty}{H_n(x)}
\end{align}

Показатели информационной эффективности:
\begin{itemize}
	\item производительность источника;
	\item информационная насыщенность источника;
	\item избыточность источника.
\end{itemize}

Производительность источника
\begin{align}
	--
\end{align}

Информационная насыщенность источника
\begin{align}
	--
\end{align}

Избыточность источника
\begin{align}
	--
\end{align}

\chapter{Информационные характеристики непрерывных сообщений}

\section{Дифференциальные энтропии случайных ДИС}

Информационная функция НИС
\begin{align}
	J_n(\cdot) = -\log_b{W_n(\cdot)} 
\end{align}

Дифференциальная энтропия
\begin{gather}
	n=1: H[\xi] = M[J(\xi)] = -M\log_b{W_1(\xi)} = -\int_{-\infty}^{+\infty} W_1(x)\log_b{W_1(x)}dx \\
	n=2: H[\xi_1, \xi_2] =  M[J(\xi_1, \xi_2)] = -M\log_b{W_1(\xi_1, \xi_2)} =\nonumber\\= -\int_{-\infty}^{+\infty}\int_{-\infty}^{+\infty} W_2(x_1, x_2)\log_b{W_1(x_2, x_1)}dx_1dx_2\nonumber
\end{gather}
Условная дифференциальная энтропия

Взаимная дифференциальная энтропия

Примеры одномерных НИС и их энтропии:
\begin{itemize}
	\item НИС с равномерной ФПВ:
\begin{gather}
	W_1(x) = 
	\begin{cases}
		\frac{1}{L_x}, & x_{\min} < x < x_{\max},\\
		0, & \text{иначе},
	\end{cases}, \\
	H[\xi] = \log_b{L_x}, m_x = \frac{x_{\min} + x_{\max}}{2}, D_\xi = \frac{L_x^2}{12}, L_x = x_{\max} - x_{\min}; \nonumber
\end{gather}
	\item НИС с гауссовской ФПВ:
\begin{gather}
	W_1(x) = \frac{1}{\sqrt{2\pi D_\xi}}e^{-\frac{x^2}{2D_\xi}}, \\
	H[\xi] = \log_b{\sqrt{2\pi D_\xi e}}, M[\xi] = 0, M[\xi^2] = D_\xi; \nonumber
\end{gather}
	\item НИС с экспоненциальной ФПВ:
\begin{gather}
	W_1(x) = 
	\begin{cases}
		\lambda e^{-\lambda x}, & 0 \leq x < +\infty,\\
		0, & \text{иначе},
	\end{cases}, \\
	H[\xi] = \log_b{\frac{e}{\lambda}}, m_\xi = \frac{1}{\lambda}, D_\xi = \frac{1}{\lambda^2}.\nonumber
\end{gather}
\end{itemize}

\section{Экстремальные свойства энтропии}

В отличии от ДИС, у непрерывных случайных сообщений диапазон возможных значений неограничен. Поэтому экстремальная постановка задачи достижения максимального значения дифференциальной энтропии обычно рассматривается для двух наиболее важных случаев:
\begin{itemize}
	\item ограниченная шкала изменения значений случайного НИС в заданном 
	диапазоне $\Delta_x = [x_{\min}, x_{\max}]$ длиной $L_x$:
	\begin{align}
		H[\xi] \leq H^{(a)}_{\max} = n \cdot \log_b{\sqrt{2\pi D_\xi e}}
	\end{align}
	\item неограниченная шкала возможных значений НИС, но при ограничениях на 
	математическое ожидание и дисперсию:
	\begin{align}
		H[\xi] \leq H^{(b)}_{\max} = n \cdot \log_b{\sqrt{2\pi D_\xi e}},
	\end{align}
	где предполагается, что $\xi$ -- произвольное n-мерное случайное сообщение с зависимыми или независимыми элементами, полученными при дискретизации НИС.
\end{itemize}

\chapter{Обозначения и сокращения}

\begin{table}[!htbp]
	\caption{Обозначения}
	\begin{tabularx}{\textwidth}{|p{6em}|p{5em}|p{7em}|X|}
		\hline
		\textbf{Обозначение} & \textbf{Единицы} & \textbf{Термин} & \textbf{Определение} \\\hline
		$s(t)$ & В & Сигнал & Функция во времени \\\hline
		$P(jf)$ & В·с & Комплексный спектр & Преобразование Фурье \\\hline
		$G(f)$ & В²·с²/Гц & Спектральная плотность энергии & $|P(jf)|^2$ \\\hline
		$B(\tau)$ & В²·с & Корреляционная функция & Автокорреляция \\\hline
		$F_e$ & Гц & Эффективная ширина спектра & -- \\\hline
		$\tau_e$ & с & Эффективная длительность & -- \\\hline
		$B_0$ & -- & База сигнала & $F_e\tau_e$ \\\hline
		$n_{\mathrm{dim}}$ & -- & Размерность сигнала & $2\lfloor B_0\rfloor$ \\\hline
		$I$ & бит & Информация & -- \\\hline
		$H$ & бит & Энтропия & -- \\\hline
	\end{tabularx}
\end{table}
\clearpage

\begin{table}[!htbp]
	\caption{Сокращения}
	\begin{tabularx}{\textwidth}{|p{2.5em}|p{5em}|p{10em}|X|}
		\hline
		\textbf{Глава} & \textbf{Сокращение} & \textbf{Расшифровка} & \textbf{Описание} \\\hline
		1 & АЧХ & Амплитудно-частотная характеристика & -- \\\hline
		1 & НЧ & Низкочастотный & -- \\\hline
		1 & ЧО & Частотно-ограниченный & -- \\\hline
		5 & ДИС & Дискретное информационное сообщение & -- \\\hline
		5 & ДИИ & Дискретный источник информации & -- \\\hline
		5 & ДИИ-БП & Дискретный источник информации без памяти & Независимые символы \\\hline
		5 & ДИИ-С & Стационарный дискретный источник информации & -- \\\hline
		5 & ДРВ & Дискретное распределение вероятностей & -- \\\hline
		5 & ФПВ & Функция плотности вероятности & -- \\\hline
		6 & НИС & Непрерывное информационное сообщение & -- \\\hline
		6 & ИКМ & Импульсно-кодовая модуляция & -- \\\hline
	\end{tabularx}
\end{table}